\providecommand{\Lcal}{\mathcal{L}}
\providecommand{\rH}{\mathrm{H}}
\providecommand{\Fcal}{\mathcal{F}}

\chapter{Analytic Iwasawa Comparison}
\label{chap:analytic-comparison}

% archon:covers MR4372220OnTheAnticyclotomicIwasawaTheoryOfRationalEllipticCurvesAtEisensteinPrimes/Analytic.lean

Throughout this chapter, \(E/\mathbb{Q}\) is an elliptic curve of conductor \(N\), \(p\nmid 2N\) is a prime of good reduction for \(E\), and \(K\) is an imaginary quadratic field satisfying the Heegner, split, and discriminant hypotheses of Castella--Grossi--Lee--Skinner: in particular \(p = v\bar{v}\) splits in \(K\). We assume \(E[p]^{ss} \simeq \mathbb{F}_p(\phi)\oplus\mathbb{F}_p(\psi)\) as \(G_{\mathbb{Q}}\)-modules, with \(\phi,\psi\) labeled so that \(p\nmid\mathrm{cond}(\phi)\) and \(\phi\neq\mathds{1}\). We fix an integral ideal \(\mathfrak{N}\subset\mathcal{O}_K\) with \(\mathcal{O}_K/\mathfrak{N}\simeq\mathbb{Z}/N\mathbb{Z}\). The Iwasawa algebra \(\Lambda\) and the modules \(M_\theta\), \(M_E\), and the Euler factors \(\mathcal{P}_w(\cdot)\) are as in \cref{def:iwasawa-algebra}.

\section{\texorpdfstring{\(p\)}{p}-adic \texorpdfstring{\(L\)}{L}-functions}

\subsection{The Bertolini--Darmon--Prasanna \texorpdfstring{\(p\)}{p}-adic \texorpdfstring{\(L\)}{L}-function}

\begin{theorem}[BDP anticyclotomic \(p\)-adic \(L\)-function]
  \label{thm:BDP}
  \lean{MR4372220.Analytic.bdpLFunction}
  \uses{def:iwasawa-algebra}
  % SOURCE: [Castella--Grossi--Lee--Skinner, MR4372220, §Analytic side, p-adic L-functions, Thm.~\ref{thm:BDP}] (read from references/MR4372220-anticyclotomic-iwasawa-source/Eisenstein.tex)
  % SOURCE QUOTE: "There exists an element $\Lcal_E\in\Lambda^{\rm ur}$ characterized by the following interpolation property: For every character $\xi$ of $\Gamma$ crystalline at both $v$ and $\bar{v}$ and corresponding to a Hecke character of $K$ of infinity type $(n,-n)$ with $n\in\Z_{>0}$ and $n\equiv 0\pmod{p-1}$, we have $\mathcal{L}_E(\xi)=\frac{\Omega_p^{4n}}{\Omega_\infty^{4n}}\cdot\frac{\Gamma(n)\Gamma(n+1)\xi(\mathfrak{N}^{-1})}{4(2\pi)^{2n+1}\sqrt{D_K}^{2n-1}}\cdot\bigl(1-a_p\xi(\bar{v})p^{-1}+\xi(\bar{v})^2p^{-1}\bigr)^2\cdot L(f/K,\xi,1)$, where $\Omega_p$ and $\Omega_\infty$ are CM periods attached to $K$ as in \cite[\S{2.5}]{cas-hsieh1}."
  \textit{Source: Castella--Grossi--Lee--Skinner, MR4372220, §Analytic side, \S{}\(p\)-adic \(L\)-functions, Theorem~BDP.}

  There exists an element \(\Lcal_E \in \Lambda^{\mathrm{ur}}\) characterized by the following interpolation property. For every character \(\xi\) of \(\Gamma\) crystalline at both \(v\) and \(\bar{v}\) and corresponding to a Hecke character of \(K\) of infinity type \((n,-n)\) with \(n\in\mathbb{Z}_{>0}\) and \(n\equiv 0\pmod{p-1}\), we have
  \[
    \Lcal_E(\xi)
    = \frac{\Omega_p^{4n}}{\Omega_\infty^{4n}}
      \cdot\frac{\Gamma(n)\Gamma(n+1)\,\xi(\mathfrak{N}^{-1})}{4\,(2\pi)^{2n+1}\sqrt{D_K}^{2n-1}}
      \cdot\bigl(1 - a_p\xi(\bar{v})p^{-1} + \xi(\bar{v})^2 p^{-1}\bigr)^2
      \cdot L(f/K,\xi,1),
  \]
  where \(\Omega_p\) and \(\Omega_\infty\) are the CM periods attached to \(K\), \(f\in S_2(\Gamma_0(N))\) is the newform associated with \(E\), and the Euler factor \(\bigl(1-a_p\xi(\bar{v})p^{-1}+\xi(\bar{v})^2p^{-1}\bigr)^2\) is nonzero under our hypotheses.
\end{theorem}

% SOURCE QUOTE PROOF: "This was originally constructed in \cite{BDP} as a continuous function of $\xi$, and later explicitly constructed as a measure in \cite{cas-hsieh1} (following the approach in \cite{brakocevic}). Since this refined construction will be important for our purposes in this section, we recall some of the details."
\begin{proof}
  The element \(\Lcal_E\) was originally constructed in Bertolini--Darmon--Prasanna as a continuous function of \(\xi\), and later explicitly constructed as a measure by Castella--Hsieh, following the approach of Brakočević. The construction proceeds via Serre--Tate theory: for a class \([\mathfrak{a}]\) in \(\mathrm{Pic}(\mathcal{O}_c)\) (for an auxiliary integer \(c\) prime to \(Np\)), the newform \(f\) admits a Serre--Tate expansion \(f(t_{\mathfrak{a}})\in\mathbb{Z}_p^{\mathrm{ur}}\llbracket t_{\mathfrak{a}}-1\rrbracket\), which defines a \(\mathbb{Z}_p^{\mathrm{ur}}\)-valued measure \(\mathrm{d}\mu_{f,\mathfrak{a}}\) on \(\mathbb{Z}_p\) via Mahler's theorem. Using the \(p\)-depletion \(f^\flat\) and an auxiliary anticyclotomic Hecke character \(\eta\) of infinity type \((1,-1)\) and conductor \(c\), one forms the measure \(\mathscr{L}_{v,\eta}\in\Lambda^{\mathrm{ur}}\) by integrating against \(\eta_v\)-twists. The element \(\Lcal_E\) is then defined by \(\Lcal_E(\xi) := \mathscr{L}_{v,\eta}(\eta^{-1}\xi)^2\), and the stated interpolation formula follows from the computation in Castella--Hsieh.
\end{proof}

\subsection{Katz anticyclotomic \texorpdfstring{\(p\)}{p}-adic \texorpdfstring{\(L\)}{L}-functions}

\begin{theorem}[Katz anticyclotomic \(p\)-adic \(L\)-function]
  \label{thm:Katz}
  \lean{MR4372220.Analytic.katzLFunction}
  \uses{def:iwasawa-algebra}
  % SOURCE: [Castella--Grossi--Lee--Skinner, MR4372220, §Analytic side, Katz p-adic L-functions, Thm.~\ref{thm:Katz}] (read from references/MR4372220-anticyclotomic-iwasawa-source/Eisenstein.tex)
  % SOURCE QUOTE: "There exists an element $\mathcal{L}_{\theta}\in\Lambda^{\rm ur}$ characterized by the following interpolation property: For every character $\xi$ of $\Gamma$ crystalline at both $v$ and $\bar{v}$ and corresponding to a Hecke character of $K$ of infinity type $(n,-n)$ with $n\in\Z_{>0}$ and $n\equiv 0\pmod{p-1}$, we have $\Lcal_{\theta}(\xi)=\frac{\Omega_p^{2n}}{\Omega_\infty^{2n}}\cdot 4\Gamma(n+1)\cdot\frac{(2\pi i)^{n-1}}{\sqrt{D_K}^{n-1}}\cdot\bigl(1-\theta^{-1}(p)\xi^{-1}(v)\bigr)\cdot\bigl(1-\theta(p)\xi(\bar{v})p^{-1})\bigr)\times\prod_{\ell\vert C}(1-\theta(\ell)\xi(w)\ell^{-1})\cdot L(\theta_{K}\xi\mathbf{N}_K,0)$"
  \textit{Source: Castella--Grossi--Lee--Skinner, MR4372220, §Analytic side, Katz \(p\)-adic \(L\)-functions, Theorem~\(\ref{thm:Katz}\).}

  Let \(\theta : G_{\mathbb{Q}}\to\mathbb{Z}_p^\times\) be a Dirichlet character of conductor \(C\) with \(C\mid N\), and let \(\mathfrak{C}\mid\mathfrak{N}\) be such that \(\mathcal{O}_K/\mathfrak{C} = \mathbb{Z}/C\mathbb{Z}\). There exists an element \(\Lcal_\theta\in\Lambda^{\mathrm{ur}}\) characterized by the following interpolation property. For every character \(\xi\) of \(\Gamma\) crystalline at both \(v\) and \(\bar{v}\) and corresponding to a Hecke character of \(K\) of infinity type \((n,-n)\) with \(n\in\mathbb{Z}_{>0}\) and \(n\equiv 0\pmod{p-1}\), we have
  \begin{align*}
    \Lcal_\theta(\xi)
    = \frac{\Omega_p^{2n}}{\Omega_\infty^{2n}}\cdot 4\Gamma(n+1)\cdot\frac{(2\pi i)^{n-1}}{\sqrt{D_K}^{n-1}}
      &\cdot\bigl(1-\theta^{-1}(p)\xi^{-1}(v)\bigr)\cdot\bigl(1-\theta(p)\xi(\bar{v})p^{-1}\bigr)\\
      &\quad\times\prod_{\ell\mid C}(1-\theta(\ell)\xi(w)\ell^{-1})\cdot L(\theta_K\xi\mathbf{N}_K,0),
  \end{align*}
  where \(\Omega_p, \Omega_\infty\) are as in \cref{thm:BDP}, \(\theta_K\) is the base change of \(\theta\) to \(K\), \(\mathbf{N}_K\) is the norm character of \(K\), and for each \(\ell\mid C\) we take \(w\mid\ell\) with \(w\mid\mathfrak{C}\). Moreover, since all primes dividing the conductor of \(\theta\) split in \(K\), Hida's theorem (\cite{hidamu=0}) gives \(\mu(\Lcal_\theta) = 0\).
\end{theorem}

% SOURCE QUOTE PROOF: "The character $\theta$ (viewed as a character of $K$) defines a projection $\pi_\theta:\bZ_p^{\rm ur}\llbracket{\rm Gal}(K(\mathfrak{C}p^\infty)/K)\rrbracket\rightarrow\Lambda^{\rm ur}$, where $K(\mathfrak{C}p^\infty)$ is the ray class field of $K$ of conductor $\mathfrak{C}p^\infty$ (this projection is just $g\mapsto \theta(g)[g]$ for $g\in {\rm Gal}(K(\mathfrak{C}p^\infty)/K$ and $[g]$ the image of $g$ in $\Gamma$). The element $\Lcal_{\theta}$ is then obtained by applying $\pi_\theta$ to the Katz $p$-adic $L$-function described in \cite[Thm.~27]{kriz}, setting $\chi^{-1}=\theta_K\xi\mathbf{N}_K$."
\begin{proof}
  The character \(\theta\), viewed as a character of \(K\), defines a projection
  \[
    \pi_\theta : \mathbb{Z}_p^{\mathrm{ur}}\llbracket\mathrm{Gal}(K(\mathfrak{C}p^\infty)/K)\rrbracket \longrightarrow \Lambda^{\mathrm{ur}},
  \]
  where \(K(\mathfrak{C}p^\infty)\) is the ray class field of \(K\) of conductor \(\mathfrak{C}p^\infty\), given by \(g\mapsto\theta(g)[g]\) for \(g\in\mathrm{Gal}(K(\mathfrak{C}p^\infty)/K)\) and \([g]\) its image in \(\Gamma\). The element \(\Lcal_\theta\) is obtained by applying \(\pi_\theta\) to the Katz \(p\)-adic \(L\)-function of Katz--Hida--Tilouine with the substitution \(\chi^{-1} = \theta_K\xi\mathbf{N}_K\). The vanishing \(\mu(\Lcal_\theta) = 0\) follows from Hida's result, since the conductors of \(\phi\) and \(\psi\) divide \(N\) and all prime divisors of \(N\) split in \(K\).
\end{proof}

\section{Comparison II: Analytic Iwasawa invariants}

\begin{theorem}[Kriz Eisenstein congruence]
  \label{thm:kriz}
  \lean{MR4372220.Analytic.krizCongruence}
  \uses{thm:BDP, thm:Katz}
  % SOURCE: [Castella--Grossi--Lee--Skinner, MR4372220, §Analytic side, Comparison II, Thm.~\ref{thm:kriz}] (read from references/MR4372220-anticyclotomic-iwasawa-source/Eisenstein.tex)
  % SOURCE QUOTE: "$\mathcal{L}_E\equiv(\mathcal{E}_{\phi,\psi}^\iota)^2\cdot(\mathcal{L}_\phi)^2\pmod{p\Lambda^{\rm ur}}$, where $\mathcal{E}_{\phi,\psi}=\prod_{\ell\vert N_0N_-}\mathcal{P}_w(\phi)\cdot\prod_{\ell\vert N_0N_+}\mathcal{P}_w(\psi)$"
  \textit{Source: Castella--Grossi--Lee--Skinner, MR4372220, §Analytic side, Comparison II: Analytic Iwasawa invariants, Theorem~\(\ref{thm:kriz}\).}

  Assume \(E[p]^{ss}\simeq\mathbb{F}_p(\phi)\oplus\mathbb{F}_p(\psi)\) as \(G_{\mathbb{Q}}\)-modules, with \(\phi,\psi\) labeled so that \(p\nmid\mathrm{cond}(\phi)\) and \(\phi\neq\mathds{1}\). Let \(N_0\) be the square-full part of \(N\) (so \(N/N_0\) is square-free), and let \(\lambda^\iota\) denote the image of \(\lambda\in\Lambda\) under the involution \(\gamma\mapsto\gamma^{-1}\). There is a factorization \(N/N_0 = N_+ N_-\) with
  \[
    \begin{cases}
      a_\ell \equiv \phi(\ell)\pmod{p} & \text{if } \ell\mid N_+,\\
      a_\ell \equiv \psi(\ell)\pmod{p} & \text{if } \ell\mid N_-,\\
      a_\ell \equiv 0\pmod{p}          & \text{if } \ell\mid N_0,
    \end{cases}
  \]
  such that the following congruence holds:
  \[
    \Lcal_E \equiv (\mathcal{E}_{\phi,\psi}^\iota)^2 \cdot (\Lcal_\phi)^2 \pmod{p\Lambda^{\mathrm{ur}}},
  \]
  where
  \[
    \mathcal{E}_{\phi,\psi}
    = \prod_{\ell\mid N_0 N_-}\mathcal{P}_w(\phi)
      \cdot\prod_{\ell\mid N_0 N_+}\mathcal{P}_w(\psi),
  \]
  and for each \(\ell\mid N\) we take the prime \(w\mid\ell\) with \(w\mid\mathfrak{N}\).
\end{theorem}

% SOURCE QUOTE PROOF: "By Theorem~34 and Remark~32 in \cite{kriz}, our hypothesis on $E[p]$ implies that there is a congruence $f\equiv G^{}\pmod{p}$, where $G^{}$ is a certain weight two Eisenstein series (denoted $E_{2}^{\phi,\phi^{-1},(N)}$ in \emph{loc.\,cit.}). Viewed as a $p$-adic modular form, $G$ defines $\Z_p^{\rm ur}$-valued measures $\mu_{G,\mathfrak{a}}$ on $\Z_p$ by the rule (\ref{eq:def-meas})."
\begin{proof}
  \uses{thm:BDP, thm:Katz}
  By Theorem~34 and Remark~32 of Kriz, the hypothesis on \(E[p]\) implies a congruence of \(p\)-adic modular forms \(f\equiv G\pmod{p}\), where \(G\) is a weight-two Eisenstein series of type \(E_2^{\phi,\phi^{-1},(N)}\). Viewing \(G\) as a \(p\)-adic modular form, its Serre--Tate expansions define \(\mathbb{Z}_p^{\mathrm{ur}}\)-valued measures \(\mu_{G,\mathfrak{a}}\) on \(\mathbb{Z}_p\) by the same rule as in the proof of \cref{thm:BDP}. The corresponding measure element \(\Lcal_G \in \Lambda^{\mathrm{ur}}\) is defined by \(\Lcal_G(\xi) := \mathscr{L}_{v,\eta}(G, \eta^{-1}\xi)\). The interpolation computation of Kriz (Prop.~37 of \emph{op.\,cit.}, with \(\chi^{-1} = \xi\mathbf{N}_K\), \(\psi_1 = \phi\), \(\psi_2 = \phi^{-1} = \psi\omega^{-1}\)) then gives \(\Lcal_G = \mathcal{E}_{\phi,\psi}^\iota\cdot\Lcal_\phi\), using that \(\mathcal{E}_{\phi,\psi}(\xi^{-1}) = \Xi_{\xi^{-1}\mathbf{N}_K^{-1}}(\phi,\psi\omega^{-1}, N_+, N_-, N_0)\). On the other hand, the congruence \(f\equiv G\pmod{p}\) implies \(\Lcal_E \equiv (\Lcal_G)^2 \pmod{p\Lambda^{\mathrm{ur}}}\). Combining these two identities yields the stated congruence.
\end{proof}

\begin{corollary}[mu/lambda formula for \(\Lcal_E\)]
  \label{cor:Kriz}
  \lean{MR4372220.Analytic.krizLambda}
  \uses{thm:kriz, thm:Katz, def:iwasawa-algebra}
  % SOURCE: [Castella--Grossi--Lee--Skinner, MR4372220, §Analytic side, Comparison II, Thm.~\ref{cor:Kriz}] (read from references/MR4372220-anticyclotomic-iwasawa-source/Eisenstein.tex)
  % SOURCE QUOTE: "Assume that $E[p]^{ss}=\mathbb{F}_p(\phi)\oplus\mathbb{F}_p(\psi)$ as $G_\bQ$-modules, with the characters $\phi, \psi$ labeled so that $p\nmid\mathrm{cond}(\phi)$, and suppose $\phi\neq\mathds{1}$. Then $\mu(\Lcal_E)=0$ and $\lambda(\Lcal_E)=\lambda(\Lcal_\phi)+\lambda(\Lcal_\psi)+\sum_{w\in S}\bigl\{\lambda(\mathcal{P}_w(\phi))+\lambda(\mathcal{P}_w(\psi))-\lambda(\mathcal{P}_w(E))\bigr\}$."
  \textit{Source: Castella--Grossi--Lee--Skinner, MR4372220, §Analytic side, Comparison II, Theorem~\(\ref{cor:Kriz}\).}

  Under the hypotheses of \cref{thm:kriz}, we have \(\mu(\Lcal_E) = 0\) and
  \[
    \lambda(\Lcal_E)
    = \lambda(\Lcal_\phi) + \lambda(\Lcal_\psi)
      + \sum_{w\in S}\bigl\{\lambda(\mathcal{P}_w(\phi)) + \lambda(\mathcal{P}_w(\psi)) - \lambda(\mathcal{P}_w(E))\bigr\},
  \]
  where \(S\) is the set of primes \(w\nmid p\) in \(\Sigma\) split in \(K\).
\end{corollary}

% SOURCE QUOTE PROOF: "Since $K$ satisfies (\ref{eq:intro-Heeg}) and (\ref{eq:intro-spl}), the conductors of both $\phi$ and $\psi$ are only divisible by primes split in $K$, and hence the vanishing of $\mu(\mathcal{L}_E)$ follows immediately from the congruence of Theorem~\ref{thm:kriz} and Hida's result \cite{hidamu=0} (note that the factors $\mathcal{P}_w(\phi)$ and $\mathcal{P}_w(\psi)$ also have vanishing $\mu$-invariant, since again the primes $w$ are split in $K$)."
\begin{proof}
  \uses{thm:kriz}
  Since \(K\) satisfies the Heegner and split hypotheses, the conductors of \(\phi\) and \(\psi\) are only divisible by primes split in \(K\). Therefore \(\mu(\Lcal_E) = 0\) follows from the congruence of \cref{thm:kriz} and Hida's theorem \(\mu = 0\) for Katz \(p\)-adic \(L\)-functions (note that the Euler factor terms \(\mathcal{P}_w(\phi)\) and \(\mathcal{P}_w(\psi)\) also have vanishing \(\mu\)-invariant, since the primes \(w\) are split in \(K\)).

  For the \(\lambda\)-invariant formula, the involution \(\gamma\mapsto\gamma^{-1}\) of \(\Lambda\) preserves \(\lambda\)-invariants, so
  \[
    \lambda(\mathcal{P}_w(\theta)^2) = \lambda(\mathcal{P}_w(\theta)) + \lambda(\mathcal{P}_{\bar{w}}(\theta)),
  \]
  using that complex conjugation acts as inversion on \(\Gamma\). For the term \(\mathcal{E}_{\phi,\psi}\) from \cref{thm:kriz} we therefore get
  \[
    \lambda((\mathcal{E}_{\phi,\psi}^\iota)^2)
    = \sum_{w\mid N_0 N_-}\lambda(\mathcal{P}_w(\phi)) + \sum_{w\mid N_0 N_+}\lambda(\mathcal{P}_w(\psi)),
  \]
  where \(w\) runs over all divisors. Using the congruence relations in \cref{thm:kriz} (in particular \(a_\ell\equiv 0\pmod{p}\) for \(\ell\mid N_0\)), this rewrites as
  \[
    \lambda((\mathcal{E}_{\phi,\psi}^\iota)^2)
    = \sum_{w\in S}\bigl\{\lambda(\mathcal{P}_w(\phi)) + \lambda(\mathcal{P}_w(\psi)) - \lambda(\mathcal{P}_w(E))\bigr\}.
  \]
  Since \(\psi = \phi^{-1}\omega\), the functional equation for the Katz \(p\)-adic \(L\)-function gives \(\lambda(\Lcal_\psi) = \lambda(\Lcal_\phi)\). The result now follows from \cref{thm:kriz} combined with these two identities.
\end{proof}

\begin{theorem}[Analytic-algebraic lambda comparison]
  \label{mulambda}
  \lean{MR4372220.Analytic.mulambda}
  \uses{cor:Kriz, MAINalgside}
  % SOURCE: [Castella--Grossi--Lee--Skinner, MR4372220, §Analytic side, Comparison II, Thm.~\ref{mulambda}] (read from references/MR4372220-anticyclotomic-iwasawa-source/Eisenstein.tex)
  % SOURCE QUOTE: "Assume that $E[p]^{ss}=\mathbb{F}_p(\phi)\oplus\mathbb{F}_p(\psi)$ with $\phi\vert_{G_p}\neq\mathds{1},\omega$. Then $\mu(\mathcal{L}_E)=\mu(\mathfrak{X}_E)=0$ and $\lambda(\Lcal_E)=\lambda(\mathfrak{X}_E)$."
  \textit{Source: Castella--Grossi--Lee--Skinner, MR4372220, §Analytic side, Comparison II, Theorem~\(\ref{mulambda}\).}

  Assume \(E[p]^{ss} = \mathbb{F}_p(\phi)\oplus\mathbb{F}_p(\psi)\) with \(\phi|_{G_p}\neq\mathds{1},\omega\). Then
  \[
    \mu(\Lcal_E) = \mu(\mathfrak{X}_E) = 0
    \quad\text{and}\quad
    \lambda(\Lcal_E) = \lambda(\mathfrak{X}_E).
  \]
\end{theorem}

% SOURCE QUOTE PROOF: "The vanishing of $\mu(\mathfrak{X}_E)$ (resp. $\mu(\mathcal{L}_E)$) has been shown in Proposition~\ref{propmodp} (resp. Theorem~\ref{cor:Kriz}). On the other hand, Iwasawa's main conjecture for $K$ (a theorem of Rubin \cite{rubinmainconj}) yields in particular the equalities $\lambda(\Lcal_\phi)=\lambda(\mathfrak{X}_\phi)$ and $\lambda(\Lcal_\psi)=\lambda(\mathfrak{X}_\psi)$. The combination of Theorem~\ref{MAINalgside} and Theorem~\ref{cor:Kriz} therefore yields the result."
\begin{proof}
  \uses{cor:Kriz, MAINalgside}
  The vanishing \(\mu(\mathfrak{X}_E) = 0\) is the algebraic conclusion of \cref{MAINalgside}, and \(\mu(\Lcal_E) = 0\) is \cref{cor:Kriz}. For the \(\lambda\)-equality, Rubin's proof of Iwasawa's main conjecture for \(K\) gives in particular \(\lambda(\Lcal_\phi) = \lambda(\mathfrak{X}_\phi)\) and \(\lambda(\Lcal_\psi) = \lambda(\mathfrak{X}_\psi)\). Substituting these equalities into the formula of \cref{cor:Kriz} and comparing with the algebraic comparison formula of \cref{MAINalgside} (which expresses \(\lambda(\mathfrak{X}_E)\) in terms of \(\lambda(\mathfrak{X}_\phi)\), \(\lambda(\mathfrak{X}_\psi)\), and the local Euler factors) yields \(\lambda(\Lcal_E) = \lambda(\mathfrak{X}_E)\).
\end{proof}
